%% kapa: exam=1 round=20 year=2009 period=2 form=A subject=accounting subject_ko=회계학 seq=29 number=29
%% source: Q-Net 감정평가사 기출문제 (cst00309) artlSeq=5024178 file=extracted/1008/5024178/2009년 제20회 감정평가사 1차시험 기출문제/2교시.hwp
%% conversion: kordoc 4.24.51 -> tools/gen1.py (auto); review=proofread-sonnet; flags=none
\begin{question}{29}
(주)한국은 구 기계장치를 처분하고 새 기계장치를 20×6년 1월 1일에 ￦300,000을 지급하고 구입하였다. 새 기계장치에 대하여 운임 및 설치비가 ￦100,000 발생하였으며, 사용가능한 상태로 만들기 위한 시운전비용이 ￦30,000 발생하였다. 새 기계장치의 내용연수는 8년이며, 잔존가치는 ￦30,000이다. 새 기계장치는 정액법에 의해 감가상각한다. 20×8년 1월 1일에 새 기계장치에 취득원가 ￦60,000인 부속장치를 부착하였으며, 이 부속장치는 새 기계장치의 생산능력을 증대시켰다. 부속장치의 부착이 기계장치의 감가상각방법, 내용연수 및 잔존가치에 미치는 영향은 없다. 20×8년 4월 1일에 새 기계장치의 유지․보수를 위해 ￦30,000을 지출하였다. 20×8년 12월 31일 현재 이 기계장치의 장부금액은 얼마인가? (단, 감가상각은 월할 계산한다.)
\begin{choices}
  \item ￦282,500
  \item ￦320,000
  \item ￦330,000
  \item ￦340,000
  \item ￦350,000
\end{choices}
\end{question}
\answer{3}
